\chapter{Giriş}
\label{ch:giris}

\section{Konu}

Ticari taksi hizmeti, kentsel ulaşımın temel bileşenlerinden biridir;
buna karşın sürücü ve yolcunun birbirini önceden tanımadığı, kapalı
mekânda ve genellikle gece saatlerinde gerçekleşen bu etkileşim, güvenlik
açısından kırılgan bir noktadır. Sürücüye yönelik saldırı, gasp ve
benzeri olaylar; hem can güvenliği hem de mesleki sürdürülebilirlik
bakımından önemli risk faktörleridir. Mevcut çözümlerin çoğu (panik
butonu, GPS izleme, dijital takograf) ya \emph{olay sonrası} bilgi
sağlamakta ya da sürücünün aktif müdahalesini gerektirmektedir.

Bu çalışmada, sürücünün dikkatini ve sürüş güvenliğini bozmadan,
\textbf{tek bir butona basarak} arka koltuktaki yolcunun aranan kişiler
veri tabanı ile eşleştirilmesini sağlayan, düşük maliyetli ve modüler
bir gömülü sistem prototipi tasarlanmıştır. Sistem; görüntü yakalama,
ön~işleme, bulut tabanlı yüz tanıma, sürücüye görsel/işitsel geri
bildirim ve otomatik 155 çağrısı bileşenlerini uçtan uca entegre eder.

\section{Problem}

Taksi sürücüleri, yolcu kaynaklı tehdit veya saldırı vakalarına karşı
açık bir çalışma ortamına sahiptir. Bu riskleri azaltmak için iki temel
ihtiyaç tanımlanmıştır:

\begin{enumerate}[label=(\roman*),itemsep=2pt]
    \item Yolcu kimliğinin caydırıcılık oluşturacak biçimde
          doğrulanması ve kayıt altına alınması,
    \item Bir olay durumunda hızlı müdahaleyi sağlayacak, sürücünün
          aktif eylemine bağımlı olmayan bir bildirim mekanizmasının
          bulunması.
\end{enumerate}

Var olan çözümler bu iki ihtiyacı tek bir gömülü sistemde, düşük
donanım maliyetiyle ve sürücünün dikkatini dağıtmayacak biçimde bir
araya getirememektedir. Bu eksiklik, mevcut tezin çıkış noktasıdır.

\section{Amaç ve Kapsam}

Tezin amacı, yukarıda tanımlanan iki ihtiyacı karşılayan bir prototip
sistem geliştirmek ve uçtan uca işlevselliğini deneysel olarak
göstermektir. Proje kapsamı aşağıdaki sınırlarla belirlenmiştir:

\begin{itemize}[itemsep=2pt]
    \item Tek aracı kapsayan, sergi ve laboratuvar düzeneğinde test
          edilebilen prototip seviyesinde bir sistem geliştirilecektir.
    \item Kamera ve yüz tanıma; gerçek bir taksi içi montaj senaryosu
          yerine, tavan döşemesi/arka koltuk karşısı ideal koşullar
          baz alınarak değerlendirilecektir.
    \item Acil çağrı yalnızca \textbf{155 Polis İmdat} numarasına
          yönlendirilecek; SMS/e-posta benzeri ek bildirim katmanları
          mimaride yer ayrılmış olsa da bu tezin değerlendirme
          kapsamı dışındadır.
    \item Yolcu görüntüleri sunucu tarafında diske yazılmaz; yalnızca
          işleme süresince RAM'de tutulur. Kalıcı saklama politikası
          KVKK çerçevesinde tartışılır.
\end{itemize}

\section{Hedeflenen Başarı Ölçütleri}
\label{sec:hedefler}

Gelişme raporunda \cite{ekin_gelisme_2026} tanımlanan hedef metriklerden
yola çıkılarak, bitirme tezi kapsamında aşağıdaki başarı ölçütleri
belirlenmiştir:

\begin{itemize}[itemsep=2pt]
    \item \textbf{Uçtan uca tanıma gecikmesi:} ESP32-CAM TARA komutu
          alındıktan, STM32'ye sonucun ulaştığı ana kadar geçen süre
          $\leq$~7~s (10 karelik öbek için, Wi-Fi $\geq$~10\,Mbps yukarı
          hızda).
    \item \textbf{Yüz tanıma doğruluğu (ideal koşullarda):}
          Eşit Hata Oranı (EER) $\leq$~\%5; FAR\,$\leq$\,\%1 noktasındaki
          FRR $\leq$~\%10. Ölçüm \cref{sec:olcum-plani}'ndeki test
          protokolüne göre yapılacaktır.
    \item \textbf{Aydınlatma ve mesafe aralığı:} 100--600~lx aydınlatma ve
          30--120~cm yolcu--kamera mesafesi aralığında kararlı çalışma.
    \item \textbf{Acil çağrı tetikleme süresi:} STM32 MATCH kararından
          telefonun 155'i çevirmeye başlamasına kadar $\leq$~3~s.
    \item \textbf{Ek donanım maliyeti:} Mevcut STM32 kartı dışında
          eklenen bileşenler için $\leq$~500~\textcent\,(TL).
\end{itemize}

\paragraph{Önerilen değişiklik.} Gelişme raporundaki ``uçtan uca
$<$~200~ms'' hedefi, tek karelik bir API sorgusunu varsayıyordu. Çoklu
kare öbek (\emph{burst}) tasarımına geçildiği için bu hedef
\emph{karelerin tamamı için} en fazla 7~s olarak revize edilmiştir
(\cref{sec:degisiklikler}). Yolcu binişi başına yalnızca bir tarama
yapıldığından, kullanıcı deneyimi açısından bu süre uygundur.

\section{Tezin Yapısı}

\Cref{ch:literatur}'de yüz tanıma yöntemleri ve gömülü sistem tabanlı
uygulamalar literatürü özetlenir. \Cref{ch:yontem} sistemin tasarım
ilkelerini ve mimari kararlarını sunar. \Cref{ch:gerceklestirme} ise
gerçekleme aşamasında alınan kararları, kullanılan donanım ve yazılım
bileşenlerini, protokol tanımlarını ve karşılaşılan teknik sorunları
ayrıntılı olarak açıklar. \Cref{ch:sonuclar}, gerçekleştirilen
ölçümleri ve elde edilen sonuçları tartışır. \Cref{ch:sonuc-oneriler}
ise çıkarımları toparlar ve gelecek çalışma önerilerini paylaşır.
